\section{From APP seeds to mixture and mass registration}
\label{sec:semillas-registro-mezcla-autonoma}
\index{seed!factorization sum--product}
\index{mixture!origin in the seeds}

The physical law receives a domain that has already preserved the incompatibility of
the two APP evaluations. If
\(\mathcal S_+\) and \(\mathcal S_\times\) denote the oriented seeds of
the additive and multiplicative leaves, then
\begin{equation}
 \Omega_{\rm seed}=\mathcal S_+\times\mathcal S_\times,
 \qquad
 |\Omega_{\rm seed}|=324^2=104,976.
 \label{eq:dominio-semillas-md-autonoma}
\end{equation}
The exact factorization of the emitter produces
\begin{equation}
 104,976
 \longrightarrow18\times26=468\ U_6
 \longrightarrow243\ w_6\subset\mathbb F_3^6,
 \qquad |\mathbb F_3^6|=729.
 \label{eq:embudo-semillas-md-autonoma}
\end{equation}
The four cardinalities count distinct objects: initial conditions of
two sheets, decimal emissions, published ternary words, and the algebraic
ambient space. In addition to its emission, each seed preserves multiplicity,
diagonal, sheet, orbit, emission signature, residues modulo three and modulo
nine, turns, sheet changes, endpoint, displacements, and cardinal
counts. Those data constitute the vocabulary of the subsequent register.

The factorization organizes six microfibers of cardinality \(1008\), six phase
Fano planes, thirty-six fibers of cardinality \(28\), and the band
\(1944=27\cdot72\), diagonalized by \((\mathbb Z/3)^3\). Upon averaging the
nonadic word of the TPK, the conditional projectors of the two sheets
define the canonical kernel
\begin{equation}
 K_{\rm cat}=\frac13(I+E_++E_\times),
 \qquad
 \operatorname{Spec}(K_{\rm cat})
 =\left\{1^{[1]},
 \left(\frac23\right)^{[42]},
 \left(\frac13\right)^{[425]}\right\}.
 \label{eq:kernel-canonico-semillas-md-autonoma}
\end{equation}
The TPK phase selects the active renewal and preserves the nine
spatial bands; the average~\eqref{eq:kernel-canonico-semillas-md-autonoma}
publishes the stochastic shadow of the complete calendar.

\subsection{Bigrading of the Seeds and \(6\) Incidence Planes}

In each of the six microfibers of cardinality \(1008\), the bipartite graph
retains \(54\) additive vertices and \(56\) multiplicative vertices. The
former decompose into \(36\) vertices of degree \(24\) and \(18\) of
degree \(8\); the latter have degree \(18\). The incidence count
agrees on both sheets:
\begin{equation}
 1008=36\cdot24+18\cdot8=56\cdot18.
 \label{eq:bidegraduacion-microfibra-masas-autonoma}
\end{equation}
This equality fixes the emitter bigrading before projecting it onto
the route families.

The six words
\begin{equation}
 001112,\quad010211,\quad100121,\quad112001,\quad121100,\quad211010
 \label{eq:seis-palabras-microfibra-masas-autonoma}
\end{equation}
form an orbit of \(D_3\simeq S_3\) and have the profile
\(432+4\cdot144=1008\). The quotient \(1008=36\cdot28\) organizes seven
blocks
\[
 Q=\{B_3,B_6,B_9,P_1,P_2,P_3,f\},
\]
where \((B_3,B_6,B_9)\) are the transverse tetrads and
\((P_1,P_2,P_3,f)\), the four lateral tetrads. For one of the six
words \(a\), let
\[
 \mathcal R_a
 :=\{(p,t)\in\mathcal S_+\times\mathcal S_\times:
       U_6(p,t)\text{ publishes the word }a\}
\]
the set of its one thousand eight routes, and set
\(R_9:=\mathbb Z/9\mathbb Z\). Once the oriented chart is fixed, the phase
projection has type
\begin{equation}
 q_\beta:\mathcal R_a\longrightarrow\binom{R_9}{2},
 \qquad
 q_\beta(p,t)
 =\{x-5\beta(d_t),x+5\beta(d_t)\}\subset\mathbb Z/9\mathbb Z,
 \qquad x=5(i_p-j_p),
 \label{eq:proyeccion-fano-semillas-masas-autonoma}
\end{equation}
with \((\beta(N),\beta(E),\beta(S),\beta(O))=(1,2,3,4)\). Its codomain is
the set of the thirty-six unordered pairs of \(R_9\). Each
\(X\in Q\) is a tetrad of orbits;
its route support and its cut over a fiber are defined by
\[
 \operatorname{Supp}_{\rm rutas}(X)
 :=\bigcup_{\mathcal O\in X}\mathcal O,
 \qquad
 X_b:=\operatorname{Supp}_{\rm rutas}(X)\cap q_\beta^{-1}(b),
 \qquad b\in\binom{R_9}{2}.
\]
In the multiplicative channel, the source observable is
\begin{equation}
 \rho:\mathcal S_\times\longrightarrow R_9,
 \qquad
 \rho(i,j,d)=(i+1)(j+1)\pmod9.
 \label{eq:residuo-producto-fano-semillas-rev3}
\end{equation}
If \(\pi_\times:\mathcal R_a\to\mathcal S_\times\),
\(\pi_\times(p,t)=t\), is the second projection, the observable acting
on routes is
\begin{equation}
 \rho_\times:=\rho\circ\pi_\times:\mathcal R_a\longrightarrow R_9,
 \qquad
 \rho_\times(p,t)=(i_t+1)(j_t+1)\pmod9.
 \label{eq:residuo-producto-rutas-fano-semillas-rev3}
\end{equation}
In what follows, \(\rho(X)\) abbreviates the multiset of values of
\(\rho_\times\) over \(X_b\), which is independent of the fiber \(b\). Its
profiles over the tetrads are
\[
 \rho(B_3)=3^4,
 \qquad
 \rho(B_6)=6^4,
 \qquad
 \rho(B_9)=0^4,
\]
and, for the three lateral tetrads,
\begin{equation}
 \begin{array}{c|c|c}
 r&\text{active signature of }P_r&R(P_r)\\
 \hline
 1&933&\{5,7\}\\
 2&393&\{4,8\}\\
 3&339&\{1,2\}.
 \end{array}
 \label{eq:perfiles-laterales-fano-semillas-rev3}
\end{equation}
Every residue of \(R(P_r)\) occurs twice among the four routes of
\((P_r)_b\).

\begin{theorem}[Incidence morphism of the seeds and six Fano planes]
\label{thm:semillas-seis-fano-fase-rev3}
For \(\delta\in R_9\), \(h\in\{3,6,0\}\) ---corresponding,
respectively, to \(B_3,B_6,B_9\)--- and \(r\in\{1,2,3\}\), let
\begin{equation}
 N_\delta(B_h,P_r)
 :=\sum_{b\in\binom{R_9}{2}}
 \#\left\{(x,y)\in(B_h)_b\times(P_r)_b:
 \rho_\times(y)-\rho_\times(x)=\delta\right\}.
 \label{eq:conteo-incidencia-fano-semillas-rev3}
\end{equation}
Then
\begin{equation}
 N_\delta(B_h,P_r)=
 \begin{cases}
 288,&h+\delta\in R(P_r),\\
 0,&h+\delta\notin R(P_r).
 \end{cases}
 \label{eq:conteo-288-cero-fano-semillas-rev3}
\end{equation}
If \(\delta\in R_9^\times\), for each \(B_h\) there exists a unique \(P_r\)
with count \(288\); therefore, one obtains a bijection
\[
 \theta_\delta:\{B_3,B_6,B_9\}
 \xrightarrow{\sim}\{P_1,P_2,P_3\}.
\]
Writing the image of \((B_3,B_6,B_9)\) by means of its subscripts, the six
bijections are
\begin{equation}
 \begin{array}{c|c@{\qquad}c|c}
 \delta&\theta_\delta&\delta&\theta_\delta\\
 \hline
 1&(2,1,3)&5&(2,3,1)\\
 2&(1,2,3)&7&(3,2,1)\\
 4&(1,3,2)&8&(3,1,2).
 \end{array}
 \label{eq:tabla-theta-delta-fano-semillas-rev3}
\end{equation}
Be \(P=\{P_1,P_2,P_3\}\) and, for \(p\in P\), be
\[
 \mu(p):=\bigl\{\{f,p\},P\setminus\{p\}\bigr\}
\]
the perfect matching associated over the four lateral tetrads.
Then
\begin{equation}
 \mathcal D_\delta
 :=\bigl\{\{B_3,B_6,B_9\}\bigr\}
 \cup
 \bigl\{\{B,q,q'\}:B\in\{B_3,B_6,B_9\},
       \{q,q'\}\in\mu(\theta_\delta(B))\bigr\}
 \label{eq:sistema-fano-semillas-rev3}
\end{equation}
is a system of Steiner \(\operatorname{S}(2,3,7)\). The six units
produce the six planes of Fano over these seven blocks that contain the
line transverse \(\{B_3,B_6,B_9\}\).
\end{theorem}

\begin{proof}
Fix a fiber \(b\). In \((B_h)_b\) there are four choices with residue \(h\). If \(h+\delta\in R(P_r)\), the profile \eqref{eq:perfiles-laterales-fano-semillas-rev3} provides exactly two choices in \((P_r)_b\) with that residue; if it does not belong to the support, it provides none. Since there are thirty-six fibers, \(N_\delta=36\cdot4\cdot2=288\) in the first case and zero in the second. The three lateral supports partition \(R_9^\times\), so a unit \(\delta\) determines a unique lateral support for each transversal block. Direct reading of those supports yields the table \eqref{eq:tabla-theta-delta-fano-semillas-rev3} and exhausts the \(3!=6\) bijections.

The transversal line covers the three pairs among \(B_3,B_6,B_9\). For a fixed \(B\), the two edges of the matching \(\mu(\theta_\delta(B))\) cover the four lateral points once. The three matchings are distinct and form the factorization of \(K_4\) into its three perfect matchings; hence every lateral pair occurs exactly once. The seven triples thus contain each of the twenty-one pairs of points once, which is property \(\operatorname{S}(2,3,7)\). Conversely, every Fano plane containing the transversal line assigns to each \(B\) a perfect matching distinct from the other two and recovers one of the preceding six bijections.
\end{proof}

The correspondence of sets \(\delta\mapsto\theta_\delta\) is neither an
action nor a homomorphism: \(R_9^\times\simeq C_6\) is abelian under
multiplication, whereas \(S_3\) is not, and the units do not form an additive
subgroup. The displacement \(\delta\) indexes six translations of the product
profiles.

The same calculation and the same table are verified in all six microfibers. Nevertheless,
the difference between support and role must be preserved: the support of
\(f\) coincides with one of the lateral supports and, therefore, the observable
\(\rho\) does not by itself recover the decomposition \(P\sqcup\{f\}\). That
distinction requires the additive signature already declared in the seed state.

\begin{statusbox}[title={Status and scope of Fano incidence}]
The \cref{thm:semillas-seis-fano-fase-rev3} is a finite theorem once
\(q_\beta\), the orientation, and the digital order have been fixed: it defines six exact
incidence gauges on blocks. It organizes phase and incidence in the seed domain;
it selects neither particle, mass, nor route. Its physical realization is performed
subsequently through the morphisms of the enriched ledger.
\end{statusbox}

The passage to physics preserves the chain
\begin{equation}
 \begin{aligned}
 \text{refined seed}
 &\longrightarrow(\text{signatures }+,\times;U_6;w_6;
                   \text{sheet memory})\\
 &\longrightarrow\text{CT108 record}
 \longrightarrow(\tau_{12},D_k,\Omega,P_6,\nu_{120},\nu_{270})\\
 &\longrightarrow\text{fiber and multisection}
 \longrightarrow(B_D,B_\Omega)
 \longrightarrow\text{CKM/PMNS transport}\\
 &\longrightarrow\text{character and towers}
 \longrightarrow\text{mass section}.
 \end{aligned}
 \label{eq:cadena-semilla-mezcla-masa-autonoma}
\end{equation}
CKM performs a change of basis of the quark operators,
\begin{equation}
 B_d^{(u)}=V_{\rm CKM}B_dV_{\rm CKM}^*,
 \qquad
 R_{\rm act}^{B_d^{(u)}}
 =V_{\rm CKM}R_{\rm act}^{B_d}V_{\rm CKM}^*,
 \label{eq:ckm-transporte-semillas-autonoma}
\end{equation}
and preserves spectrum and determinant. The mixing specifies the representation
of the operator in a flavour basis; the register preserves the eigenvalues and their
genealogy.

\subsection{Neutrino reader contained in the register}

The neutral sector has an integer extractor constructed exclusively
with record data. For each dodecaphase block \(E_m\), we define
\begin{equation}
 T(m)=\sum_{t\in E_m}(-1)^{\kappa(t)}\varepsilon(t)
       \mathbf1_{\{d(t)=9\}},
 \qquad
 s_C^{(m)}(\nu)=T(m)-T(m+6),
 \label{eq:lector-neutrino-ledger-autonoma}
\end{equation}
and
\begin{equation}
 \tau_m^\nu=\operatorname{sgn}s_C^{(m)}(\nu).
 \label{eq:trit-neutrino-ledger-autonoma}
\end{equation}
The four jump triangles produce
\begin{equation}
 \nu_{120}=\sum_{j=1}^{4}
 \operatorname{sgn}\!\left(\sum_{m\in\mathcal T_j}\tau_m^\nu\right),
 \qquad
 \nu_{270}=\sum_{m=1}^{12}\tau_m^\nu\tau_{m+3}^\nu.
 \label{eq:invariantes-neutrino-ledger-autonoma}
\end{equation}
Bidirectional inversion transforms
\(\nu_{120}\mapsto-\nu_{120}\) and preserves \(\nu_{270}\). This parity
defines the historical architecture of a reader independent of metrological
magnitudes used as a target. Its literal application
to the living record \(81\times108\), taking as
\(\varepsilon(t)\) the sole global sign bit preserved by that
projection, yields
\[
 s_C^{(m)}(\nu)=0,
 \qquad \tau_m^\nu=\nu_{120}=\nu_{270}=0
\]
in the eighty-one positions. Antipodal events have been
identified by projection of the register. The local result is preserved
as a recovered reader and, simultaneously, as a test of information loss: the
neutral extractor of
\cref{sec:extractor-neutro-z0-torsion-rev2} acts before that
identification and reads the local bidirectional bit, the sheet, torsion, and the
antipodal corona. The zero obtained after projection does not demote that
construction; it certifies exactly what information the global reader erases.
The frames
\(F_\ell,F_\nu\), the reader~\eqref{eq:invariantes-neutrino-ledger-autonoma},
and the aggregate form
\begin{equation}
 G_{ab}=\frac1{\sqrt{|L_a||N_b|}}
 \sum_{c\in L_a}\sum_{d\in N_b}K_{\rm registro}(c,d),
 \qquad U_{\rm int}=\operatorname{polar}(G),
 \label{eq:kernel-ledger-pmns-autonoma}
\end{equation}
Thus, the two already constructed ends of the PMNS bridge are fixed.

The cardinal equality between the fifty-six multiplicative seed vertices and the fifty-six CT108 families has undergone an initial control independent of metrological magnitudes used as targets. The natural Hadamard map on a single seed reaches only twenty-nine families; upon traversing the natural dihedral symmetries, signs, and translations, the maximum coverage is thirty-five. The intertwiner therefore acts on the enriched seed pair and preserves sheet, carry, and memory. The incidence matrix comparing shared emissions in that double domain must have rank fifty-six, be natural under orientation, and not factor through global counters. This criterion fixes the exact domain of the next construction.
